\section{Introduction}
\label{sec:introduction}

Clinical institutions may operate vendor-supplied language models within
retrieval and memory scaffolds without access to the model weights or the
vendor's fine-tuning pipeline. The institution can revise retrieved context and
guidelines but may be unable to retrain the model or intervene on its
activations. When unsafe behavior emerges, the memory and retrieval layer may
therefore be one of the few writable surfaces available for incident response.
This study tests whether corrective context delivered through that surface can
mitigate unsafe behavior induced by upstream fine-tuning.

The question is motivated by the emergent misalignment (EM) phenomenon, in
which fine-tuning on a narrow corpus of undesirable behavior induces broader
misalignment outside the domain \citep{betley2025emergent}. Open model organisms
make this effect reproducible, including a Qwen2.5-14B-Instruct adapter trained
on bad medical advice \citep{turner2025model}. Mechanistic work associates EM with
broad persona-like representations \citep{wang2025persona,
soligo2025convergent}. Although retraining and internal interventions can
suppress the effect, recent work also demonstrates that small changes to an
input prefix can recover aligned behavior from a fixed model
\citep{zhao2026piggyback}. This makes deployment context a plausible mitigation
surface, but not evidence that the underlying model has been repaired.

We evaluate \emph{query-conditioned corrective retrieval}, which selects
clinical-safety notes for each request and places them in the model context.
The primary comparisons contrast no intervention (C1), a fixed corrective
system prompt (C2), static corrective retrieval (C3), length- and
format-matched neutral retrieval (C5), and the un-fine-tuned base model (C6).
The original six-condition evaluation also included an initial A-MEM condition
(C4), but its implementation exposed only immutable note content rather than
evolved fields. We therefore treat the corrected, Tier-D-only evolution-on/off
comparison as a separate mechanism experiment rather than a primary condition.
Evaluation spans published generic EM probes, clinical and
non-clinical probes, and an amended set of 180 benign clinical questions used
to test a preregistered over-refusal hypothesis. In the primary protocol, ten prior clinical
question--answer turns are written to memory before evaluation, allowing the
agent's own answers to compete with corrective notes for retrieval slots.

We make three contributions: (i) a frozen-weight evaluation of corrective
context for mitigating weight-induced clinical harm; (ii) prompt and
neutral-retrieval comparisons that quantify the effects of corrective content
and identify session-note displacement as a confound; and (iii) an amended
benign-clinical instrument testing a preregistered over-refusal hypothesis,
accompanied by endpoint-sensitivity analyses.
